\begin{exercisechunk}
\item \textbf{Point corrections and the test measure.}
\label[exercise]{ex:l9-point-corrections}
\exerciseprofile{2}{1}{2}{\exproof\quad\exconstruction\quad\exinterpretation}{%
Separate exact fitting at the observations from integrated prediction error.}
Let $\mu$ be a Borel probability law on $\cX\subset\R$, let
$x_1,\ldots,x_n\in\cX$ be distinct, and let $g,f^\star$ be measurable,
square-integrable functions with specified point values.
For numbers $r_1,\ldots,r_n$, define
\[
 h(x)=g(x)+\sum_{i=1}^nr_i\I\{x=x_i\}.
\]

\begin{enumerate}[(a),beginpenalty=10000]
\item Suppose that $\mu(\{x\})=0$ for every $x\in\cX$. Prove
\[
 \|h-g\|_{L^2(\mu)}=0,
 \qquad
 \|h-f^\star\|_{L^2(\mu)}
 =\|g-f^\star\|_{L^2(\mu)}.
\]
Choose $r_i$ so that $h(x_i)=y_i$ for arbitrary prescribed values
$y_1,\ldots,y_n$.

\item Show that the conclusion can fail when the test law has a point mass.
Give an explicit
choice of $\cX,\mu,g,f^\star,x_1,r_1$ for which
\[
 \|h-f^\star\|_{L^2(\mu)}^2
 >\|g-f^\star\|_{L^2(\mu)}^2.
\]
Show that the distinction from part (a) is whether the set on which the
correction is nonzero has zero or positive $\mu$-measure, and state why this
determines whether the correction contributes to the population
$L^2(\mu)$ error.

\item \textbf{Optional.} Let $P\in\cP_2$ be arbitrary, and assume that every singleton in
$\cX$ is measurable. Let
$S_n=((X_i,Y_i))_{i=1}^n$ be an i.i.d. sample with
$\Law(X_i,Y_i)=P$. Write
$N_n(x)=\sum_{i=1}^n\I\{X_i=x\}$ and define
\[
 \widehat f_n(x)=
 \begin{cases}
  N_n(x)^{-1}\sum_{i:X_i=x}Y_i,&N_n(x)>0,\\
  \widetilde f_n(x),&N_n(x)=0,
 \end{cases}
\]
where $\widetilde f_n$ is any estimator. Prove that, for every training index
$i$ such that
$X_j=X_i$ implies $Y_j=Y_i$ for all $j\in[n]$,
\[
 \widehat f_n(X_i)=Y_i.
\]

\item \textbf{Optional.} In the setting of part (c), suppose that
\[
 \|\widetilde f_n-f_P^\star\|_{L^2(P_X)}^2\xrightarrow{\P}0.
\]
Prove
\[
 \|\widehat f_n-f_P^\star\|_{L^2(P_X)}^2\xrightarrow{\P}0
 \qquad\text{as }n\to\infty.
\]
You may use that the set
$\cD=\{x:P_X(\{x\})>0\}$ is countable, the strong law of large numbers, and
dominated convergence for sums over $\cD$.
\end{enumerate}
\begin{hints}
For (b), a one-point space is sufficient.
\end{hints}
\end{exercisechunk}
